% ══════════════════════════════════════════════════════════════════════════════
%  EXPERIMENT 6 — GROWING COPPER SULFATE CRYSTALS               (Lesson 14, Term 2 Week 5)
% ══════════════════════════════════════════════════════════════════════════════
\experimentdiv{Growing Copper Sulfate Crystals}{Handling a hazardous chemical \textbullet{} Safe disposal \textbullet{} Measuring growth over two weeks}{Lesson 14 \textbullet\ Term 2, Week 5 \textbullet\ Stretch lab}\label{sec:6}
\lessoncodes{VC2S8U05, U07 \textbullet{} VC2S8I02, I07, I08}

\vocabhead
\begin{ctab}
\begin{tabular}{|SW{4.2cm}|W{11.6cm}|}
  \hline\hd & \thead{Word} & \thead{Definition (copy from the slide)} \\ \hline
  & \scaf{\textbf{Seed crystal}} & \scaf{A small crystal that the dissolved particles build onto as the solution evaporates.} \\ \hline
  & \textbf{Hazard} & \\ \hline
  & \textbf{Risk} & \\ \hline
  & \textbf{Disposal} & \\ \hline
  & \textbf{Ethical} & \\ \hline
\end{tabular}
\end{ctab}

\begin{lipanel}
\inquiryli{Identify hazards and plan safe handling and disposal; record observations over time.}
\understandli{Explain how a crystal grows as the solvent evaporates from a saturated solution.}
\sctitle
\begin{itemize}[leftmargin=1.4em, itemsep=0pt, topsep=1pt]
  \item I can name the hazards of copper sulfate and how to control them.
  \item I can explain why copper sulfate must not go down the sink.
  \item I can measure and draw my crystal as it grows.
\end{itemize}
\end{lipanel}

\begin{skillsfocus}
\focus{PC3}{Name the hazards and control the risks.}
\focus{PC5}{Dispose of chemicals responsibly.}
\focus{PC6}{Measure the crystal every few days.}
\focus{CO2}{Draw the set-up and the crystal.}
\focus{EV5}{Suggest how to grow a bigger crystal.}
\end{skillsfocus}

% ─── TECHNIQUE ────────────────────────────────────────────────────────────────
\exppart{Technique spotlight}{Blue}

\begin{technique}{Risk, prevention and responding}{PC3}
A \textbf{hazard} is something that could cause harm. The \textbf{risk} is what could go wrong with it. For each risk, plan the \textbf{prevention} (what you do so it does not happen) and the \textbf{response} (what you do if it happens anyway).
\end{technique}

\begin{workedex}{One row of a risk assessment}
\textbf{Risk:} hot copper sulfate solution is spilled on skin.\quad \textbf{Prevention:} the teacher pours it; students keep beakers in the middle of the bench.\quad \textbf{Responding:} rinse the skin under cool running water and tell the teacher.
\end{workedex}

\begin{youdobox}[Risk assessment]
\task Complete the risk assessment for this experiment. \skr{PC3}
\begin{ctab}
\begin{tabular}{|SW{3.8cm}|W{5.0cm}|W{5.4cm}|}
  \hline\hd & \thead{Risk (what could go wrong)} & \thead{Prevention} & \thead{Responding (if it happens)} \\ \hline
  & \scaf{Copper sulfate is swallowed} & \scaf{Never taste; wash hands after handling} & \scaf{Rinse mouth with water; tell the teacher at once} \\ \hline
  & Copper sulfate gets in eyes or on skin & & \\ \hline
  & The glass beaker breaks & & \\ \hline
  & Hot solution is spilled & & \\ \hline
\end{tabular}
\end{ctab}
\end{youdobox}

\begin{technique}{Responsible disposal}{PC5}
Copper sulfate is \textbf{toxic to fish and other water life}. Water from sinks flows to the sewage plant and on to rivers and the bay, and the treatment plant does not remove all of it. Scientists have a responsibility to the environment as well as to themselves: every drop of copper sulfate solution goes into the labelled \textbf{heavy metal waste} container, never down the sink. Crystals you keep stay in a sealed, labelled bag.
\end{technique}

\begin{wedo}{Why it matters}
A student says, ``It's only a few drops, it won't matter if I tip it down the sink.'' Explain why the class still uses the waste container.
\answerbox{3.5cm}
\end{wedo}

% ─── PLAN AND SAFETY / METHOD ────────────────────────────────────────────────
\exppart{Plan, safety and method}{Teal}

\textbf{Aim:} To grow a large crystal of copper sulfate from a seed crystal.

\subsubsection*{Equipment}
100\,mL beaker \textbullet{} saturated copper sulfate solution (made hot by your teacher) \textbullet{} a small seed crystal \textbullet{} cotton thread \textbullet{} pencil or paddle-pop stick \textbullet{} paper towel and rubber band \textbullet{} ruler \textbullet{} hand lens \textbullet{} gloves and safety glasses

\begin{center}
\begin{tikzpicture}[scale=0.8]
  \eqBeaker
  \fill[black!15] (0.02,0.02) rectangle (1.78,1.8); \draw[Ink, line width=0.8pt] (0.02,1.8) -- (1.78,1.8);
  \draw[eq] (-0.5,2.75) -- (2.3,2.75);
  \draw[Ink, line width=0.8pt] (0.9,2.75) -- (0.9,1.0);
  \draw[eq, fill=black!55] (0.9,1.05) -- (1.1,0.85) -- (1.0,0.6) -- (0.8,0.6) -- (0.7,0.85) -- cycle;
  \eqlab{(2.2,2.75)}{(3.2,3.1)}{pencil across the top}
  \eqlab{(0.9,1.4)}{(3.2,2.2)}{cotton thread}
  \eqlab{(1.1,0.85)}{(3.2,1.2)}{seed crystal, off the bottom}
  \eqlab{(1.78,0.4)}{(3.2,0.4)}{saturated copper sulfate solution}
\end{tikzpicture}
\expfigcap{The seed crystal hangs in the solution without touching the glass.}
\end{center}

\begin{infobox}{Safety}
Wear safety glasses and gloves. Copper sulfate is harmful if swallowed and irritates eyes and skin: never taste it, and wash your hands after the lesson even if you wore gloves. Your teacher pours the hot solution. All leftover solution goes in the heavy metal waste container.
\end{infobox}

\subsubsection*{Method}
\begin{enumerate}[itemsep=2pt]
  \item Tie the seed crystal to one end of the thread and the other end to the pencil.
  \item Your teacher half fills your beaker with warm saturated solution.
  \item Lay the pencil across the beaker so that the crystal hangs in the middle, clear of the bottom and the sides.
  \item Cover the beaker with paper towel held on by a rubber band, to keep dust out but let water evaporate. Label it.
  \item Leave it somewhere still. Measure and draw the crystal at each lesson for two weeks.
  \item At the end, put the leftover solution in the waste container and dry the crystal on paper towel.
\end{enumerate}

% ─── RESULTS ─────────────────────────────────────────────────────────────────
\exppart{Results}{Amber}

\subsubsection*{Table 6.1: Growth of the copper sulfate crystal \skr{PC6}}
\begin{ctab}
\begin{tabular}{|SW{2.6cm}|C{2.6cm}|W{9.8cm}|}
  \hline\hd & \thead{Day} & \thead{Longest side (mm)} & \thead{Observations (colour, shape, crystals on the glass)} \\ \hline
  & \scaf{Example} & \scaf{3} & \scaf{Blue, four-sided; a crust forming at the water line} \\ \hline
  & Day 0 & & \\ \hline
  & Day 2 & & \\ \hline
  & Day 5 & & \\ \hline
  & Day 7 & & \\ \hline
  & Day 10 & & \\ \hline
  & Day 14 & & \\ \hline
\end{tabular}
\end{ctab}

\needspace{14\baselineskip}
\task Draw your crystal through the hand lens on the first and last days. Label its faces and edges, and write its length under each drawing. \skr{CO2}

\noindent
\begin{minipage}[t]{0.48\linewidth}\drawbox{6cm}{}\nameline{Day 0}\end{minipage}\hfill
\begin{minipage}[t]{0.48\linewidth}\drawbox{6cm}{}\nameline{Day 14}\end{minipage}

% ─── ANALYSIS ────────────────────────────────────────────────────────────────
\exppart{Analysis and evaluation}{Coral}

\task Explain how the crystal grew. Use the words \emph{saturated}, \emph{evaporate} and \emph{particles}.
\answerbox{5cm}

\needspace{10\baselineskip}
\task Small crystals often formed on the sides of the beaker too. Suggest why, and why they take copper sulfate away from your seed crystal.
\answerbox{4cm}

\needspace{10\baselineskip}
\task Why must the solution be saturated before the seed crystal goes in? What would happen to the seed crystal if it were not?
\answerbox{4cm}

\needspace{10\baselineskip}
\task Explain why the leftover solution went into the waste container and not down the sink. \skr{PC5}
\answerbox{4cm}

\needspace{10\baselineskip}
\task Suggest one change that could grow a larger crystal, and write it as a question you could investigate. \skr{EV5}
\answerbox{4cm}

\begin{yr8box}
\begin{keepext}
\exttask Copper sulfate crystals are blue, but if you heat them they turn white and give off steam. Suggest what the crystals contain besides copper sulfate, and how you could test your idea.
\answerbox{5cm}
\end{keepext}
\end{yr8box}

% ─── SKILL CHECK ─────────────────────────────────────────────────────────────
\exppart{Skill check}{NavyBlue}

\begin{skillcheck}
\scrow{PC3}{I can name the hazards of copper sulfate and how to control them.}
\scrow{PC5}{I disposed of the solution responsibly and can explain why.}
\scrow{PC6}{I measured and recorded my crystal at each lesson.}
\scrow{CO2}{I drew my crystal clearly and labelled it.}
\scrow{EV5}{I can suggest a question for a further investigation.}
\end{skillcheck}

\task Which skill will you focus on in Experiment 7? Explain what you will do differently.
\answerbox{5cm}
